\section*{Disclaimer}

This document was written as part of a personal research project on Physics-Informed Neural Networks (PINNs). It is primarily intended for scientists and engineers with a background in physics, mechanical engineering, applied mathematics, or numerical simulation who are interested in transitioning toward Scientific Machine Learning (SciML) and PINNs.

Readers unfamiliar with machine learning may benefit from completing an introductory machine learning course before studying some of the more advanced concepts discussed herein. The author personally completed the \href{https://www.coursera.org/specializations/machine-learning-introduction}{Machine Learning Specialization} and \href{https://www.coursera.org/specializations/deep-learning}{Deep Learning Specialization} offered by DeepLearning.AI and taught by Andrew Ng. While these courses are not specifically focused on SciML or PINNs they provide a broad introduction to the fundamental concepts of modern machine learning and offer a sufficient foundation for readers with a scientific or engineering background to approach the methodologies presented in this document.

Large Language Models (LLMs) were occasionally used to assist with writing, proofreading, and structuring this paper. However, all technical discussions, implementations, analyses, interpretations, and conclusions were reviewed and validated by the author. Unless otherwise stated, all figures presented in this report were generated by the author.

\section{Introduction}

\subsection{Artificial Intelligence and Scientific Research}

Over the last decade, Artificial Intelligence (AI) has profoundly transformed modern society. Advances in machine learning, deep learning, and more recently Large Language Models (LLMs) have led to major breakthroughs in numerous domains, including natural language processing (the automatic understanding and generation of human language), computer vision (the extraction of information from images and videos), healthcare, finance, autonomous systems, robotics, and industrial optimization. These technologies are increasingly capable of automating complex tasks, extracting information from massive datasets, and assisting human decision-making processes. While many early applications focused on classification, prediction, and pattern recognition, AI is now increasingly being used to support scientific research itself.

Scientific disciplines have traditionally relied on a combination of theoretical modeling, experimental measurements, and numerical simulations. Compared with many machine learning application domains, scientific applications often impose additional requirements such as physical consistency, interpretability, reproducibility, and rigorous validation. This naturally creates a certain degree of inertia, as new computational methodologies must undergo extensive testing before gaining widespread acceptance. Despite these challenges, researchers have increasingly embraced machine learning techniques as powerful tools for accelerating discovery, improving data analysis, enhancing predictive modeling capabilities, and complementing traditional numerical simulation methods.

\subsection{Scientific Machine Learning}

The convergence of machine learning and scientific computing has given rise to a rapidly growing field known as Scientific Machine Learning (SciML). The objective of SciML is to combine data-driven learning with physics-based modeling. Machine learning provides flexible tools for learning complex relationships directly from data, whereas physical models incorporate established scientific knowledge about the system under study. By leveraging both sources of information, SciML seeks to improve predictive accuracy, reduce data requirements, and ensure greater physical consistency of the resulting models. Several important families of SciML methods have recently emerged, including:

\begin{itemize}
\item Physics-Informed Neural Networks (PINNs),
\item Deep Operator Networks (DeepONets),
\item Neural Operators,
\item Fourier Neural Operators (FNOs),
\item Graph Neural Networks (GNNs),
\item and hybrid physics-data models.
\end{itemize}

Deep Operator Networks (DeepONets) and Neural Operators can be viewed as machine learning methods that attempt to learn how to solve a family of physical problems rather than a single instance of a problem. Their objective is to build models capable of rapidly predicting approximate solutions once they have been trained on a sufficiently large number of examples. Fourier Neural Operators (FNOs) are a particular type of Neural Operator that have recently attracted significant attention due to their ability to efficiently approximate solutions of certain classes of partial differential equations. Graph Neural Networks (GNNs) are designed to work with systems that can naturally be represented as interconnected objects. Such representations are frequently encountered in engineering applications involving complex geometries, computational meshes, molecular structures, or networked systems.

Although these approaches differ considerably in their implementation and underlying mathematical foundations, they share a common goal: integrating machine learning with scientific modeling to enhance predictive capabilities, accelerate numerical simulations, and extract meaningful information from physical systems. Current applications of SciML in physics include, but are not limited to:

\begin{itemize}
\item fluid mechanics (e.g., turbulence modeling and flow reconstruction),
\item heat transfer (e.g., thermal property estimation and surrogate modeling),
\item solid mechanics (e.g., stress and deformation prediction),
\item electromagnetics (e.g., wave propagation and antenna design),
\item climate science (e.g., reduced-order weather forecasting),
\item geophysics (e.g., subsurface property identification),
\item materials science (e.g., microstructure-property relationships),
\item plasma physics (e.g., fusion reactor modeling),
\item biomedical engineering (e.g., blood flow and medical imaging),
\item and quantum systems (e.g., Schrödinger equation solvers and quantum control).
\end{itemize}

\subsection{Physics-Informed Neural Networks (PINNs)}

Among the various SciML methodologies, Physics-Informed Neural Networks (PINNs) have emerged as one of the most actively studied approaches since the seminal work of Raissi et al (\textit{Journal of Computational Physics}, 2019). The key idea behind PINNs is relatively simple: instead of learning solely from data, the neural network is constrained to satisfy the governing physical equations during training.

In a traditional supervised machine learning problem, a neural network is provided with examples of known inputs and outputs. The objective is to progressively adjust the network parameters so that its predictions become as close as possible to the observed outputs. For example, a neural network may be trained to predict the temperature at different positions and times within a metal bar using measurements collected during a heat diffusion experiment. By comparing its predictions to the measured temperatures, the network gradually learns the relationship between the inputs (position and time) and the desired output (temperature). To quantify the discrepancy between predictions and observations, a numerical quantity called the \textit{loss function} is introduced. The loss can be viewed as a measure of prediction error that the optimization algorithm attempts to minimize during training. The larger the prediction error, the larger the loss.

PINNs extend this idea by incorporating physical knowledge directly into the loss function. In addition to matching the available data, the network is also penalized whenever its predictions violate the governing differential equations. This is made possible through \textit{automatic differentiation}, a computational technique that allows derivatives of the neural network output to be computed efficiently and accurately. These derivatives are then used to evaluate how well the network satisfies the underlying physical equations. In the heat diffusion example, the network is therefore encouraged not only to reproduce the measured temperatures, but also to satisfy the heat equation throughout the domain. As a result, PINNs generate solutions that are consistent not only with the available data but also with the governing physics.

This approach is particularly attractive for engineering applications because many practical problems involve both physical knowledge and limited experimental data. In such situations, neither purely data-driven methods nor purely physics-based approaches may be sufficient on their own. Data-driven methods often require large amounts of high-quality data to achieve accurate predictions, which may be difficult or expensive to obtain in many scientific and engineering applications. Conversely, purely physics-based models rely on accurate governing equations, boundary conditions, and material parameters, which are not always fully known in practice. By combining observational data with physical knowledge, PINNs aim to leverage the strengths of both approaches while mitigating their respective limitations.

Since their introduction, PINNs have been successfully applied to a wide variety of problems, including heat diffusion, fluid mechanics, wave propagation, elasticity, electromagnetics, inverse parameter identification, uncertainty quantification, data assimilation, multiphysics systems, and more recently quantum systems and quantum control problems. One of the most promising applications of PINNs concerns inverse problems. In many engineering situations, physical parameters cannot be measured directly and must instead be inferred from sparse and potentially noisy observations. Examples include estimating thermal diffusivities, material properties, permeability coefficients, reaction rates, or boundary conditions from experimental measurements. By simultaneously incorporating observational data and physical constraints within a unified optimization framework, PINNs offer an attractive alternative to traditional parameter estimation techniques.

\subsection{Objectives of the Present Work}

The present work focuses on PINNs applied to a simple one-dimensional heat diffusion problem. Although the one-dimensional heat equation has already been extensively studied in the PINN literature, it remains one of the few benchmark problems for which every aspect of the learning process can be investigated in a fully controlled environment. The availability of an analytical solution makes it possible to isolate and study individual mechanisms such as optimization dynamics, parameter identifiability, loss balancing, convergence behavior, and robustness to noisy data without the additional complexity introduced by more realistic applications. Furthermore, many of the difficulties encountered in large-scale PINN applications already emerge in this simplified setting, including optimization stiffness, parameter non-identifiability, sensitivity to noise, collocation point selection (points in the computational domain where the governing equations are enforced), and convergence issues.

The study considers both the forward and inverse formulations of the problem. In the forward problem, the thermal diffusivity coefficient is assumed to be known, and the objective is to reconstruct the temperature field satisfying the governing heat equation together with the prescribed initial and boundary conditions. In the inverse problem, the objective is to estimate the thermal diffusivity coefficient $\alpha$ directly from noisy temperature measurements while simultaneously reconstructing the temperature field. This task is representative of many real-world scientific and engineering applications in which physical parameters cannot be measured directly and must instead be inferred from experimental observations.

Particular attention is devoted to understanding why PINNs may succeed or fail even on apparently simple benchmark problems. Rather than viewing the one-dimensional heat equation as a trivial example, this work treats it as a controlled laboratory for studying parameter identifiability, optimization dynamics, noise robustness, loss balancing strategies, sensitivity-based weighting approaches, hard versus soft constraint formulations, and Galerkin-inspired techniques designed to improve inverse parameter reconstruction.

The remainder of this document is organized as follows. Section 2 presents the PINN architecture, training strategies, initialization procedures, normalization techniques, and reproducibility considerations. Section 3 introduces the different machine learning frameworks investigated throughout this work. Section 4 describes the one-dimensional heat diffusion problem and derives its analytical solution. Section 5 investigates the forward PINN formulation, while Section 6 focuses on the inverse problem and the various strategies developed to improve parameter identification. Finally, the results and conclusions are presented together with perspectives for future developments.

Beyond the one-dimensional benchmark considered here, the long-term objective is to develop the methodological foundations required to tackle more realistic scientific and engineering problems for which analytical solutions are unavailable, such as systems governed by nonlinear partial differential equations or involving complex geometries. In such situations, validation must rely on experimental measurements, high-fidelity numerical simulations, or both. For the author, whose background lies in engineering physics, experimental research, and numerical simulation, this project also represents a transition toward the field of Scientific Machine Learning and a broader effort to understand how machine learning techniques can complement established scientific methodologies.
